\chapter{Дополнительные материалы вычислительного эксперимента}
\label{app:additional_materials}

\section{Параметры эксперимента и структура данных}
\label{app:params}

\begin{table}[H]
\centering
\caption{Структура выходных файлов вычислительного комплекса}
\label{tab:output_files}
\begin{tabularx}{\textwidth}{|l|X|}
\hline
\textbf{Файл / Директория} & \textbf{Содержание} \\ \hline
\texttt{suite\_result.json} & Полный результат серии: конфигурация сценариев, число репликаций, точечные оценки метрик, доверительные интервалы, стационарные распределения \\ \hline
\texttt{aggregated\_summary.csv} & Сводная таблица по всем сценариям в удобном табличном виде \\ \hline
\texttt{metric\_summaries.csv} & Таблица статистических характеристик по отдельным целевым метрикам \\ \hline
\texttt{run\_summaries.csv} & Детализированные результаты независимых репликаций (Монте-Карло) \\ \hline
\texttt{suite\_summary.txt} & Текстовая сводка параметров запущенной серии \\ \hline
\texttt{plots/} & Автоматически сгенерированные графики для базовых экспериментов \\ \hline
\texttt{plots\_reworked3/} & Графики совместного анализа чувствительности (тепловые карты, срезы) \\ \hline
\end{tabularx}
\end{table}

\textbf{Пример конфигурационного файла (фрагмент \texttt{suite\_result.json}):}
\begin{verbatim}
{
  "suite_name": "baseline_gpu",
  "replications": 40,
  "max_time": 50000.0,
  "warmup_time": 5000.0,
  "servers_n": 96,
  "total_resource_r": 590,
  "arrival_rate_levels": [90.0, 104.0, 116.0],
  "service_speed_levels": [1.2],
  "resource_values": 
  [2, 4, 8, 12, 16],
  "resource_probabilities": 
  [0.25, 0.25, 0.3, 0.15, 0.05],
  "workload_family_basic": 
  [
    "deterministic", 
    "exponential", 
    "erlang_4", 
    "hyperexp_heavy"
    ],
  "arrival_process_family": 
  [
    "poisson", 
    "erlang_2", 
    "erlang_4", 
    "hyperexp_2"
    ]
}
\end{verbatim}

\clearpage
\section{Дополнительные графики чувствительности к workload}
\label{app:workload_plots}

На рисунке~\ref{fig:app_workload_boxplots} показано, что дисперсия результатов Монте-Карло для различных метрик на уровне нагрузки $\lambda=90$ сопоставима с изменениями от типа распределения времени обслуживания, что подтверждает вывод об инвариантности.

\begin{figure}[H]
    \centering
    \begin{subfigure}{0.48\textwidth}
        \includegraphics[width=\textwidth]{workload_sensitivity_boxplot_loss_probability_sigma-1p2_lambda-90.png}
        \caption{Вероятность потери}
    \end{subfigure}
    \hfill
    \begin{subfigure}{0.48\textwidth}
        \includegraphics[width=\textwidth]{workload_sensitivity_boxplot_mean_sojourn_time_sigma-1p2_lambda-90.png}
        \caption{Среднее время пребывания}
    \end{subfigure}
    \caption{Разброс целевых метрик по репликациям в зависимости от workload ($\lambda = 90$)}
    \label{fig:app_workload_boxplots}
\end{figure}

\begin{figure}[H]
    \centering
    \begin{subfigure}{0.48\textwidth}
        \includegraphics[width=\textwidth]{workload_sensitivity_stationary_distribution_sigma-1p2_lambda-90.png}
        \caption{Общий вид $\pi(k)$}
    \end{subfigure}
    \hfill
    \begin{subfigure}{0.48\textwidth}
        \includegraphics[width=\textwidth]{workload_sensitivity_stationary_tail_sigma-1p2_lambda-90.png}
        \caption{Хвостовые вероятности}
    \end{subfigure}
    \caption{Стационарное распределение числа заявок в системе при изменении workload ($\lambda = 90$)}
    \label{fig:app_workload_stationary}
\end{figure}


\clearpage
\section{Дополнительные графики чувствительности к входящему потоку}
\label{app:arrival_plots}

На рисунке~\ref{fig:app_arrival_boxplots} проиллюстрировано статистически значимое влияние типа входящего потока на метрики при невысокой нагрузке ($\lambda=90$). 

\begin{figure}[H]
    \centering
    \begin{subfigure}{0.48\textwidth}
        \includegraphics[width=\textwidth]{arrival_sensitivity_boxplot_loss_probability_sigma-1p2_lambda-90.png}
        \caption{Вероятность потери}
    \end{subfigure}
    \hfill
    \begin{subfigure}{0.48\textwidth}
        \includegraphics[width=\textwidth]{arrival_sensitivity_boxplot_mean_sojourn_time_sigma-1p2_lambda-90.png}
        \caption{Среднее время пребывания}
    \end{subfigure}
    \caption{Разброс целевых метрик по репликациям в зависимости от типа входящего потока ($\lambda = 90$)}
    \label{fig:app_arrival_boxplots}
\end{figure}

\begin{figure}[H]
    \centering
    \begin{subfigure}{0.48\textwidth}
        \includegraphics[width=\textwidth]{arrival_sensitivity_rejection_breakdown_sigma-1p2_lambda-90.png}
        \caption{Структура отказов (серверы/ресурс)}
    \end{subfigure}
    \hfill
    \begin{subfigure}{0.48\textwidth}
        \includegraphics[width=\textwidth]{arrival_sensitivity_stationary_tail_sigma-1p2_lambda-90.png}
        \caption{Хвост стационарного распределения}
    \end{subfigure}
    \caption{Дополнительные метрики чувствительности к входящему потоку ($\lambda = 90$)}
    \label{fig:app_arrival_misc}
\end{figure}


\clearpage
\section{Дополнительные графики совместной чувствительности}
\label{app:combined_plots}

В данном разделе приведены дополнительные тепловые карты и графики относительных отклонений (delta) для тех метрик эффективности, которые не вошли в основной текст третьей главы.

\begin{figure}[H]
    \centering
    \begin{subfigure}{0.48\textwidth}
        \includegraphics[width=\textwidth]{01_heatmap_workload_arrival_mean_num_jobs_sigma-1p2.png}
        \caption{Среднее число заявок в системе}
    \end{subfigure}
    \hfill
    \begin{subfigure}{0.48\textwidth}
        \includegraphics[width=\textwidth]{06_joint_delta_vs_det_poisson_mean_num_jobs_sigma-1p2.png}
        \caption{Смещение числа заявок от базлайна (п.п.)}
    \end{subfigure}
    \caption{Совместное влияние распределений на загруженность системы}
    \label{fig:app_combined_num_jobs}
\end{figure}

На рисунке~\ref{fig:app_combined_compare} показаны срезы чувствительности при фиксации одного из распределений. Левый график демонстрирует изменение метрик при регулярном обслуживании, а правый --- при регулярном входящем потоке заявок.

\begin{figure}[H]
    \centering
    \begin{subfigure}{0.48\textwidth}
        \includegraphics[width=\textwidth]{02_compare_arrivals_fixed_workload_loss_probability_sigma-1p2.png}
        \caption{Сравнение потоков (фиксированный workload)}
    \end{subfigure}
    \hfill
    \begin{subfigure}{0.48\textwidth}
        \includegraphics[width=\textwidth]{03_compare_workloads_fixed_arrival_loss_probability_sigma-1p2.png}
        \caption{Сравнение workload (фиксированный поток)}
    \end{subfigure}
    \caption{Срезы совместной чувствительности для вероятности потери}
    \label{fig:app_combined_compare}
\end{figure}

\begin{figure}[H]
    \centering
    \begin{subfigure}{0.48\textwidth}
        \includegraphics[width=\textwidth]{09_stationary_tail_by_workload_lambda-90_sigma-1p2.png}
        \caption{$\lambda = 90$}
    \end{subfigure}
    \hfill
    \begin{subfigure}{0.48\textwidth}
        \includegraphics[width=\textwidth]{09_stationary_tail_by_workload_lambda-104_sigma-1p2.png}
        \caption{$\lambda = 104$}
    \end{subfigure}
    \caption{Хвосты стационарных распределений при различных нагрузках в комбинированных сценариях}
    \label{fig:app_combined_tails}
\end{figure}